\section*{Bab \ref{ch:g11:prob} --- Peluang dan Peubah Acak}

\begin{solution}{exo:g11:prob:1}
\omterm{def:g11:prob:model}{Model seragam} pada $52$ kartu. Hati: $\frac{13}{52} = \frac14$. Raja:
$\frac{4}{52} = \frac{1}{13}$. Hati atau raja:
$\frac{13 + 4 - 1}{52} = \frac{16}{52} = \frac{4}{13}$ (raja hati
terhitung sekali). Bukan keduanya: $1 - \frac{4}{13} = \frac{9}{13}$.
\end{solution}

\begin{solution}{exo:g11:prob:2}
$S = 6$: pasangan $(1,5), (2,4), (3,3), (4,2), (5,1)$, sehingga
$\P(S = 6) = \frac{5}{36}$.

$S \geq 10$: tiga pasangan untuk $10$, dua untuk $11$, satu untuk $12$:
$\P(S \geq 10) = \frac{6}{36} = \frac16$.

$S$ genap: jumlah $2, 4, 6, 8, 10, 12$ mempunyai $1 + 3 + 5 + 5 + 3 + 1 = 18$
pasangan: peluangnya $\frac{18}{36} = \frac12$.
\end{solution}

\begin{solution}{exo:g11:prob:3}
Peluangnya berjumlah $1$: $p = 1 - 0.3 - 0.2 - 0.4 = 0.1$. Lalu
\[
\E(X) = 0.3 \times (-1) + 0.2 \times 0 + 0.4 \times 2 + 0.1 \times 5
= -0.3 + 0.8 + 0.5 = 1 .
\]
\end{solution}

\begin{solution}{exo:g11:prob:4}
$\E(Y) = 2 \times 3 - 5 = 1$; $\V(Y) = 2^2 \times 4 = 16$;
$\sigma(Y) = 4$.
\end{solution}

\begin{solution}{exo:g11:prob:5}
Perolehannya $0$, $2$, $5$ euro dengan peluang $0.5$, $0.3$,
$0.2$; setelah dikurangi taruhan $2$ euro, $G$ mengambil nilai $-2, 0, 3$:
\begin{center}
\begin{tabular}{c|ccc}
$g$ & $-2$ & $0$ & $3$\\
\hline
$\P(G=g)$ & $0.5$ & $0.3$ & $0.2$
\end{tabular}
\end{center}
$\E(G) = -1 + 0 + 0.6 = -0.4$: pemainnya rugi $40$ sen per permainan secara
rata-rata --- jadi merugikan.
\end{solution}

\begin{solution}{exo:g11:prob:6}
Mengambil secara berurutan: $\P(X = 2) = \frac35 \times \frac24 = \frac{3}{10}$
(merah lalu merah); $\P(X = 0) = \frac25 \times \frac14 = \frac{1}{10}$
(biru lalu biru); sehingga
$\P(X = 1) = 1 - \frac{3}{10} - \frac{1}{10} = \frac{6}{10}$. Lalu
\[
\E(X) = 0 \times \frac1{10} + 1 \times \frac6{10} + 2 \times \frac3{10}
= \frac{12}{10} = 1.2 ,
\]
yang sesuai dengan gerak hati: $\frac35$ dari kedua pengambilan itu merah secara rata-rata.
\end{solution}

\begin{solution}{exo:g11:prob:7}
$\E(G^2) = 0.5 \times 4 + 0.3 \times 0 + 0.2 \times 9 = 3.8$, sehingga menurut
rumus jalan pintasnya
$\V(G) = 3.8 - (-0.4)^2 = 3.8 - 0.16 = 3.64$ dan
$\sigma(G) \approx 1.9$ euro.
\end{solution}

\begin{solution}{exo:g11:prob:8}
Manfaatnya adalah $80$ dikurangi klaimnya:
\begin{center}
\begin{tabular}{c|ccc}
$b$ & $80$ & $-920$ & $-4920$\\
\hline
$\P(B=b)$ & $0.94$ & $0.05$ & $0.01$
\end{tabular}
\end{center}
$\E(B) = 75.2 - 46 - 49.2 = -20$: pada harga ini perusahaannya \emph{rugi}
$20$ euro per polis secara rata-rata. Preminya terlalu rendah untuk risiko yang
ditanggung (preminya harus melampaui $100$ euro agar harapan manfaatnya
positif).
\end{solution}

\begin{solution}{exo:g11:prob:9}
Nilai acaknya adalah $+1$ dengan peluang $\frac14$ dan $x$ dengan
peluang $\frac34$:
\[
\E = \frac14 + \frac{3x}{4} = 0 \iff x = -\frac13 .
\]
Dengan denda $-\frac13$, menerka secara acak tidak membawa apa-apa secara rata-rata.
\end{solution}

\begin{solution}{exo:g11:prob:10}
Dua koin: keduanya sama (GG atau AA) dengan peluang $\frac24 \times 2 =
\frac12$. Keuntungan pemainnya $+1$ atau $-1$ dengan peluang yang sama:
jadi \omterm{def:g11:prob:expectation}{nilai harapannya} $0$ untuk kedua pemain --- permainannya adil.

Tiga koin: ketiganya sama (GGG atau AAA) dengan peluang $\frac{2}{8} =
\frac14$. Pemainnya menerima $2$ dengan peluang $\frac14$ dan membayar $1$
dengan peluang $\frac34$:
$\E = \frac24 - \frac34 = -\frac14$. Permainan itu menguntungkan Ani sebesar
$25$ sen per ronde; permainannya akan adil bila ia membayar $3$ euro untuk
ketiganya sama.
\end{solution}

\begin{solution}{exo:g11:prob:11}
\emph{1.} Penyelenggaranya menerima $1000 \times 5 = 5000$ euro dan membayarkan
$2000 + 5 \times 200 + 20 \times 25 = 3500$ euro seluruhnya: jadi labanya
adalah \emph{konstanta} $1500$ (tanpa keacakan sama sekali begitu semua
karcisnya terjual). Keuntungan $G$ seorang pemain (hadiah dikurangi karcis $5$ euro):
\begin{center}
\begin{tabular}{c|cccc}
$g$ & $1995$ & $195$ & $20$ & $-5$\\
\hline\rule{0pt}{11pt}%
$\P(G=g)$ & $\dfrac{1}{1000}$ & $\dfrac{5}{1000}$ & $\dfrac{20}{1000}$
& $\dfrac{974}{1000}$
\end{tabular}
\end{center}

\emph{2.}
\[
\E(G) = \frac{1995 + 5 \times 195 + 20 \times 20 - 974 \times 5}{1000}
= \frac{1995 + 975 + 400 - 4870}{1000} = -1.5 .
\]
Setiap pemain berharap rugi $1.50$ euro; atas $1000$ pemain jumlahnya
$1500$ euro --- persis sama dengan laba penyelenggaranya. Bukunya seimbang:
apa yang rugi bagi pemain secara rata-rata persis sama dengan apa yang
dipungut penyelenggaranya.
\end{solution}

\begin{solution}{pb:g11:prob:1}
\textbf{1.} $G$ mengambil $-2, -1, 0, 1, 2, 3$, masing-masing dengan
peluang $\frac16$: $\E(G) = \frac{-2 - 1 + 0 + 1 + 2 +
3}{6} = \frac12$ euro.

\textbf{2.} $\E(G^2) = \frac{4 + 1 + 0 + 1 + 4 + 9}{6} =
\frac{19}{6}$; $V(G) = \frac{19}{6} - \frac14 = \frac{35}{12}
\approx 2.92$; $\sigma(G) \approx 1.71$.

\textbf{3.} $\E(2G - 1) = 2 \times \frac12 - 1 = 0$;
$V(2G - 1) = 4\,V(G) = \frac{35}{3}$.

\textbf{4.} $\E(\text{jumlahnya}) = 7$: jadi harga adilnya $7$ euro. Pada $8$,
harapan kerugiannya $1$ euro per permainan.

\textbf{5.} Premi adilnya: $0.05 \times 800 = 40$ euro.
Dengan menarik $60$: harapan labanya $20$ euro per polis.

\textbf{6.} Membagi $5:3$ mengganjar \emph{masa lalu} --- padahal
taruhannya menjadi milik siapa pun yang memenangkan \emph{masa depan}, dan B
masih mungkin menang. Memberikan semuanya kepada A memperlakukan kemenangan
yang mungkin sebagai kemenangan yang pasti --- padahal peluang kecil B tetap
ada harganya.

\textbf{7.} Ronde yang dimainkan: $3$; hasilnya (A menang atau B menang per
ronde): AAA, AAB, ABA, ABB, BAA, BAB, BBA, BBB. Pemain B merebut
rangkaiannya hanya dengan memenangkan ketiganya: yaitu BBB saja. A memenangkan
rangkaiannya pada $7$ dari $8$ kasus: jadi pembagian adilnya
$64 \times \frac78 = 56$ pistole untuk A, dan $8$ untuk B.

\textbf{8.} Memainkan ronde tambahan yang tidak berguna itu tidak mengubah
kemenangan siapa pun: begitu A memperoleh kemenangan keenamnya, ronde
berikutnya tinggal upacara. Namun mematok cakrawalanya tepat pada $3$ ronde
membuat semua $8$ hasilnya sama-sama mungkin --- pembilangan yang jujur menuntut
cerita yang sama panjangnya, dan memanjangkannya tidak berongkos apa pun.

\textbf{9.} Pada $5$--$5$: satu ronde penentu, $32$ untuk masing-masing. Pada
$5$--$4$: A memperoleh $\frac{64 + 32}{2} = 48$. Pada $5$--$3$: ronde
berikutnya membawa ke $64$ (bila A memenangkannya) atau ke kedudukan $5$--$4$
yang bernilai $48$: jadi bagian A adalah $\frac{64 + 48}{2} = 56$ ---
yaitu $56 : 8$ milik Fermat, yang diturunkan ulang dengan berjalan mundur.

\textbf{10.} A butuh $2$, B butuh $3$: mainkan $4$ ronde
khayalan ($16$ hasil). B merebut rangkaiannya dengan $3$ atau $4$
kemenangan B: yaitu $4 + 1 = 5$ hasil. Pembagiannya:
$B$: $64 \times \frac{5}{16} = 20$; $A$: $44$.

\textbf{11.} Nilai kini yang adil bagi imbalan masa depan yang tak pasti
adalah \omterm{def:g11:prob:expectation}{nilai harapannya} --- yaitu rata-rata berbobot peluang atas
apa yang mungkin datang. Asuransi (pertanyaan 5) dan keuangan (memberi harga
pada imbalan yang tak pasti) adalah kalimat itu yang diubah menjadi industri;
demikian pula kasino, dari sisi meja yang lain.

\textbf{12.} Taruhannya selalu berakhir di suatu tempat: kedua bagiannya
adalah \omterm{def:g11:prob:rv}{peubah acak} yang berjumlah konstanta $64$, sehingga menurut
kelinearan $\E_A + \E_B = 64$ --- jadi bagian yang adil menghabiskan taruhannya,
pada kedudukan mana pun.

\textbf{13.} Kelinearan: $\E(X + Y) = 3.5 + 3.5 = 7$, tanpa tabel
--- benar bahkan untuk dadu yang direkatkan. \omterm{def:g11:prob:variance}{Ragam} satu dadu:
$\E(X^2) - 3.5^2 = \frac{91}{6} - \frac{49}{4} = \frac{35}{12}$;
untuk dadu yang saling bebas
$V(X + Y) = \frac{35}{12} + \frac{35}{12} = \frac{35}{6}$.

\textbf{14.} Setiap tamu tertentu memperoleh topinya sendiri dengan
peluang $\frac1n$ (topinya satu di antara $n$). Setelah dijumlahkan
atas $n$ tamu:
$\E(X) = n \times \frac1n = 1$ --- jadi rata-rata satu tamu yang beruntung,
entah di pesta berisi dua topi entah dua ribu topi.

\textbf{15.} $n = 2$: penugasannya: keduanya tepat ($X = 2$) atau
tertukar ($X = 0$): $\E = 1$. $n = 3$: keenam penugasannya memberi
$X = 3, 1, 1, 1, 0, 0$: $\E = \frac{6}{6} = 1$.

\textbf{16.} Perhitungannya menjumlahkan \omterm{def:g11:prob:expectation}{nilai harapan} $n$
penunjuk perorangan (``tamu $i$ memperoleh topinya''), dan
kelinearan menjumlahkan \omterm{def:g11:prob:expectation}{nilai harapan} \emph{tanpa syarat} ---
kebergantungan akan merusak sebuah hasil kali atau sebuah \omterm{def:g11:prob:variance}{ragam}, tetapi tidak
pernah merusak jumlah \omterm{def:g11:prob:expectation}{nilai harapan}. Itulah kekuatan super kelinearan: tanpa
pertanyaan tentang kebebasan sama sekali.

\textbf{17.} Kedua keuntungannya: $\E = 0.5 \times 10 - 5 = 0$ dan
$0.005 \times 1000 - 5 = 0$. \omterm{def:g11:prob:variance}{Ragamnya} (bagi keuntungannya):
A: $\E(G^2) = 0.5 \times 25 + 0.5 \times 25 = 25$;
B: $0.005 \times 995^2 + 0.995 \times 25 \approx 4\,975$:
jadi $\sigma_A = 5$, $\sigma_B \approx 70.5$. Orang membeli B: yang
mereka beli bukan \omterm{def:g11:prob:expectation}{nilai harapan} (nol pada keduanya) melainkan
\emph{sebarannya} --- yaitu peluang kecil untuk ekor kanan yang mengubah
hidup.

\textbf{18.} Angka pertama muncul pada lemparan ke-$n$ dengan peluang
$2^{-n}$ dan membayar $2^n$: jadi setiap $n$ menyumbang
$2^{-n} \times 2^n = 1$ pada \omterm{def:g11:prob:expectation}{nilai harapannya}, yang karena itu bernilai
$1 + 1 + 1 + \dots$: takhingga. Namun tidak ada orang waras yang membayar
$1\,000$ untuk bermain sekali: dengan peluang $\frac{15}{16}$
pembayarannya paling besar $16$. \omterm{def:g11:prob:expectation}{Nilai harapan} meringkas jangka panjang
taruhan yang \emph{dapat diulang}; permainan sekali seumur hidup dengan ekor
yang liar tidak diberi harga oleh rata-ratanya.

\textbf{19.} \omterm{def:g11:prob:expectation}{Nilai harapannya}: berasuransi $-20$; tanpa asuransi
$-0.01 \times 1000 = -10$. Strategi telanjang itu menang secara rata-rata
--- bagi orang yang sanggup menyerap kerugiannya dan mengulang taruhannya
berkali-kali. Sebuah keluarga yang tidak dapat bertahan dari $-1\,000$ tidak
berada dalam usaha merata-ratakan: keluarga itu membayar $10$ lebih banyak
untuk menghapus \omterm{def:g11:prob:variance}{ragamnya}. Penjamin merata-ratakan; rumah tangga membeli kepastian.

\textbf{20.} Lahir untuk membagi taruhan di antara dua penjudi
yang tak sabar; diberdayai oleh kelinearan, yang menjumlahkan harapan tanpa
meminta izin kepada kebergantungan; buta terhadap risiko, yang justru harus
diukur oleh \omterm{def:g11:prob:variance}{ragam} --- undian, permainan Petersburg, polis keluarga tadi;
lalu dinobatkan tahun berikutnya oleh hukum bilangan besar, yang menjelaskan
mengapa rata-ratanya selalu, pada akhirnya, tertagih.

\end{solution}
